\documentclass[10pt,conference]{IEEEtran}

\usepackage{cite}
\usepackage{amsmath}
\usepackage{amssymb}
\usepackage{graphicx}
\usepackage{textcomp}
\usepackage{booktabs}
\usepackage{array}
\usepackage{url}

\begin{document}

\title{LexGraph: A Neuro-Symbolic Pipeline for Proof-Carrying\\Legal Reasoning over Indian Land-Dispute Jurisprudence}

\author{
\IEEEauthorblockN{[Member 1], [Member 2], [Member 3], [Member 4]}
\IEEEauthorblockA{School of Computing and Augmented Intelligence\\
Arizona State University\\
CSE 573: Semantic Web Mining --- Fall 2026\\
Instructor: Prof. Hasan Davulcu \quad TA: Anshul Trivedi}
}

\maketitle

\begin{abstract}
Large language models used to summarize case law or answer legal questions routinely
fabricate citations, and are poor at recognizing their own errors. We present LexGraph,
a small-scale pipeline that treats this as an extraction-and-verification problem rather
than a generation problem: an ontology-constrained ``grounding'' step lifts sentences of
a judgment into typed legal objects (claims, issues, cited authorities, holdings) instead
of letting a model free-generate an analysis; a signed citation graph, built from two
complementary sign-extraction sources spanning 1950--2025, records whether a later case
relied on, distinguished, or overruled an earlier one; and a deterministic verifier checks
every proposed authority against real outbound citation edges before an IRAC-style
(Issue/Rule/Application/Conclusion) answer is assembled, refusing to certify what it
cannot verify. We evaluate LexGraph against a closed-book LLM baseline and a dense
retrieval-augmented generation (RAG) baseline on 48 Indian Supreme Court land-dispute
judgments (six per decade, 1950s--2020s), scoring all three systems with the identical
deterministic verifier. LexGraph achieves the highest citation precision (60\% of its
case-shaped citations verify against ground truth, versus 52\% for RAG and 6.9\% for the
closed-book baseline, which reproduces the legal-hallucination failure mode reported in
prior work on our own corpus). However, LexGraph recovers only 1.4\% of the real citation
edges present in these judgments, against 9.9\% for plain RAG, because its candidate
sentence selection --- the first 30 keyword-matched sentences per judgment --- is a
coverage bottleneck, not a verification failure. We report this result along with its
diagnosis: recall concentrates almost entirely in judgments whose citations are announced
in a HEADNOTE section, a pre-2000 convention that vanishes from the corpus after 2000.
This paper documents the current state of the pipeline, the evaluation methodology used
to compare it fairly against baselines, and the specific next step --- retrieval-fed
candidate selection ahead of ontology grounding --- required to close the recall gap
without giving up the precision advantage of verification.
\end{abstract}

\begin{IEEEkeywords}
legal information extraction, ontology-directed extraction, citation verification, large
language model hallucination, signed graphs, retrieval-augmented generation, legal AI
\end{IEEEkeywords}

\section{Introduction}

Large language models are increasingly used to summarize case law, draft legal arguments,
and answer legal questions, yet they remain probabilistic text generators applied to a
domain where correctness is not optional. Dahl et al.~\cite{dahl2024} found that
general-purpose LLMs hallucinate legal facts or authorities in the majority of tested
queries, and that models are poor at recognizing their own errors. In a legal setting, a
fabricated citation, or a real citation used to support a proposition it does not actually
stand for, is not a cosmetic defect: it can misdirect a self-represented litigant,
contaminate a filing, or make a governmental decision effectively unreviewable. A system
that cannot show which rule it applied, which authority it relied on, and why the record
supports that authority cannot be meaningfully contested, which undermines due process
wherever AI-assisted legal reasoning is deployed.

This project builds a small-scale, working version of a proof-carrying legal reasoning
pipeline. Instead of asking an LLM to answer a legal question directly, we force it to
first construct an explicit, checkable structure --- the issue, the rule, the cited
authority, the supporting facts, and the outcome --- and we verify that structure against
a ground-truth citation graph before any natural-language answer is produced. The target
domain is property and land-dispute litigation before the Indian Supreme Court, using a
corpus of judgments and citation metadata made available for this course project. This
scope mirrors an ongoing ``contestable legal AI'' research direction in the instructor's
lab, deliberately narrowed here to one jurisdiction, one case type, and a sampled
evaluation set so that a complete, demonstrable pipeline is achievable within a single
semester.

The contributions of this paper are:
\begin{enumerate}
\item An ontology-directed extraction pipeline (GROUND/LEARN/EXTRACT) instantiated on a
real legal corpus, with a hand-authored ontology, 24 gold examples, and rules learned
from them rather than hand-written.
\item A merged signed citation graph covering 1950--2025 that combines a
free, high-precision regex extractor for a pre-2000 law-reporting convention with an
LLM-based classifier for post-2000 judgments, where that convention does not exist.
\item A deterministic verifier and evaluation harness that scores a closed-book LLM
baseline, a dense RAG baseline, and our pipeline with the exact same citation-matching
logic, making the three systems directly comparable rather than comparable in name only.
\item An honest empirical finding, with root-cause diagnosis: plain RAG currently
recovers more real citations than our pipeline does, because of a specific, identified,
fixable bottleneck in candidate sentence selection --- not because verification is
unhelpful. We report this rather than only reporting the metrics that favor our approach.
\end{enumerate}

This paper reflects the state of the project as of September 18, 2026: the data
pipeline, signed graph, extraction pipeline, verifier, and first cross-system evaluation
are complete; a live demo dashboard is not yet built and is discussed as future work in
Section~\ref{sec:future}.

\section{Related Work}

Legal knowledge representation has a long history of separating legal concepts from the
logic of their application. LKIF-Core~\cite{hoekstra2007} provides a modular OWL ontology
of basic legal concepts (norms, acts, roles, time), and OASIS LegalRuleML~\cite{legalruleml}
extends this with a machine-readable syntax for obligations, permissions, and defeasible
rules. These frameworks show that legal text can be lifted into structured, queryable
representations, but they were built by hand and do not scale to large modern corpora.

Sunny and Sivan-Sevilla~\cite{sunny2026} address this scaling problem directly: they use
an LLM to translate statutory eligibility rules into a formal ontology, then use
satisfiability-based (SMT) verification to check whether real welfare-benefit
explanations are legally adequate. Their central finding --- that formal verification
catches statutory violations even in explanations that look reasonable to a human reader
--- is strong evidence that neuro-symbolic verification adds something that
formal-logic-free LLM prompting cannot. LexGraph applies the same ``neural extraction plus
symbolic verification'' division of labor to a different artifact (judicial citation
reasoning instead of benefits notices) and a different verification target: does a cited
authority actually exist and get relied upon, rather than does an explanation satisfy a
statute.

Our extraction method follows an older lineage: Riloff~\cite{riloff1993} introduced
AutoSlog, which learns symbolic extraction patterns from examples using generic syntactic
templates (active-verb-plus-object, passive-verb-plus-subject, etc.). We replace Riloff's
syntactic patterns with LLM-grounded semantic patterns --- the model maps a sentence onto
a small ontology of legal roles and actions instead of a parse tree --- but keep
AutoSlog's core idea that extraction should produce inspectable, reusable rules rather
than one-off, black-box answers.

On the graph side, citation networks between judgments are naturally a signed ecosystem: a
later case can rely on, distinguish, or overrule an earlier one, and treating all edges as
equivalent ``citation'' loses the information that matters most for assessing whether a
proposition is still good law. Signed Graph Convolutional Networks~\cite{derr2018} show
that explicitly modeling positive and negative edges improves representation quality over
unsigned-graph baselines; we adapt this idea at a much lighter weight (a heuristic- or
classifier-derived sign label plus a structural score) rather than a full SGCN, given the
scope of a one-semester course project.

Finally, retrieval-augmented generation~\cite{lewis2020} is the natural
grounding-in-source-text baseline for this task, and we use it as one, rather than only
as an ablation of our own pipeline: Section~\ref{sec:results} shows it substantially
reduces, but does not eliminate, the fabrication problem that closed-book prompting
exhibits.

\section{Dataset}
\label{sec:dataset}

The provided corpus is an archive of 26,688 Indian Supreme Court judgments (1950--2025),
year-organized PDF files, plus two locally available derived assets: a citation metadata
set covering 7,493 citing cases with 144,353 total outbound citation edges, and a filtered
index of 6,954 judgments identified as land or property disputes.

The raw citation metadata mixes case-to-case citations with statute, article, and section
references under a single undifferentiated list. A rule-based classifier splits these
apart: of 144,353 edges, 32.5\% are genuine case citations, 54.4\% are statute or article
references, and 13.1\% are unresolved fragments (mostly bare sub-clause numbers trailing a
statute citation, e.g.\ ``(7)''), left for manual review rather than forced into a bucket.
Every downstream use of a citation edge in this paper goes through this classifier first;
treating the raw ``precedents'' field as case citations would silently double the apparent
citation count with statute noise.

Of the 6,954 land-dispute cases, only 2,503 have at least one outbound case citation of
their own, which is required for graph construction and for verification. We sampled 500
cases from this pool, stratified evenly across decades (roughly 62--65 per decade,
1950s--2020s), preferring more heavily-cited cases within each decade. All 500 sampled
judgment PDFs have been downloaded and verified to contain clean, fully digitized text,
with no OCR step required.

Inspecting the judgment PDFs surfaced a source of citation polarity we had not planned to
rely on: pre-2000 Indian Supreme Court judgments include a HEADNOTE section that states it
explicitly, in a recurring pattern --- a list of ``Case v. Case (Citation)'' entries
followed by a verb such as \emph{relied on}, \emph{followed}, \emph{approved} (positive),
\emph{distinguished}, \emph{disapproved}, \emph{overruled} (negative), or \emph{referred
to} (neutral). This convention is a hard boundary rather than a declining trend: 0 of 190
sampled 2000--2025 judgments contain a HEADNOTE section at all. Post-2000 sign extraction
therefore goes through an LLM classifier that locates the sentence(s) around a citation
mention in the judgment body and grounds the surrounding free text into the same
treatment classes, since the same distinguishing or overruling is expressed there with no
fixed vocabulary (e.g.\ ``We are of the view that the said judgment would not even
remotely be applicable to the facts of the present case,'' which uses none of the HEADNOTE
verbs). Both sources write to the same signed-edge format, so downstream graph
construction does not need to know which era a label came from. Merging both sources over
the full 500-case sample produced 8,467 candidate citation edges, of which 62.6\% carry a
real sign label after the merge (33.6\% positive, 26.9\% neutral, 2.1\% negative, with the
remainder left unsigned where neither extractor could establish a label).

\section{Methodology}
\label{sec:method}

The pipeline has three phases, corresponding to the three questions a contestable legal
answer must resolve: what does the text say (Phase I), how does the surrounding case law
treat it (Phase II), and can the answer be trusted before it is shown to a user (Phase
III).

\subsection{Phase I: Ontology-Directed Extraction}

Following the ontology-directed extraction (ODE) method taught in this course, extraction
is not ``ask an LLM to summarize the judgment.'' A small starter ontology is defined
first: two top-level classes, \textsc{Entity} (children: Party, Court, Judge,
LegalInstrument, PriorCase, PropertyAsset) and \textsc{Action} (children:
ProceduralAction, DisputeAction, TreatmentAction, RulingAction, DispositionAction); five
semantic roles (Agent, Patient, Authority, Forum, Time); and eight target information
extraction concepts, listed in Table~\ref{tab:concepts}, along with a terminology mapping
from keywords to classes.

\begin{table}[t]
\caption{Target information extraction concepts}
\label{tab:concepts}
\centering
\footnotesize
\begin{tabular}{@{}p{0.2\linewidth}p{0.7\linewidth}@{}}
\toprule
\textbf{Concept} & \textbf{Definition} \\
\midrule
Claim & What a party asserts is true or owed to them (ownership, possession, right to recovery). \\
Issue & The question the court frames itself as deciding. \\
RuleCited & A statute, section, or article invoked as governing law. \\
AuthorityCited & A prior case invoked as precedent. \\
AuthorityTreatment & How the citing court treated an AuthorityCited; the signed-graph edge label (Section~\ref{sec:phase2}). \\
Fact & A fact the court treats as established in the record. \\
Holding & The court's determination on an Issue. \\
Outcome & The disposition of the case (appeal allowed/dismissed, etc.). \\
\bottomrule
\end{tabular}
\end{table}

The LLM's role is narrow and constrained: given a sentence, ground it into ontology facts
(which classes and roles the words instantiate) rather than freely generate an
explanation. Given a handful of gold (sentence, concept, filler) examples --- 24 in the
current ontology, drawn directly from the downloaded judgments --- symbolic rules are
learned by generalizing across examples via their lowest common ancestor in the ontology
(e.g.\ \emph{relied on}/\emph{followed}/\emph{approved} generalize to a positive
TreatmentAction trigger). Seven rules have been learned so far; the two with the most
supporting examples are TreatmentAction+Authority~$\rightarrow$~AuthorityCited (3
examples, confidence 1.0) and RulingAction+Patient~$\rightarrow$~Holding (4 examples,
confidence 0.57). The lower confidence on the latter is not noise: ``held that $X$'' maps
to Holding, RuleCited, Issue, or AuthorityCited depending on what $X$ actually contains,
and resolving this needs content-based disambiguation (does $X$ mention a statute or a
case name), not more (class, role) pairs. This is documented as a known gap in the
ontology and is discussed further in Section~\ref{sec:limitations}.

Running GROUND on every sentence of a full judgment (100--300+ sentences) would be slow
and mostly wasted, since the current rule set only fires on a handful of trigger patterns.
Phase I therefore pre-filters to the first 30 sentences containing at least one ontology
terminology keyword before spending an LLM call on them --- a cheap candidate-generation
step in front of the expensive one. As Section~\ref{sec:results} shows, this specific
choice, not the grounding step itself, is the pipeline's current recall bottleneck.

\subsection{Phase II: Signed Legal Knowledge Graph}
\label{sec:phase2}

Extracted objects and the classified citation metadata are merged into a graph in which
cases and legal propositions are nodes. Citation edges are signed using the two
complementary sources described in Section~\ref{sec:dataset}: the HEADNOTE-based
extractor for pre-2000 cases, and the LLM-based classifier for post-2000 cases, both
writing into the same \{relied\_on, distinguished, overruled\_or\_disapproved,
referred\_to, unclear\} label space. From the resulting signed graph we compute a
structural Proof-Support Score per case --- a signed variant of PageRank/HITS, following
the modeling intuition behind signed graph learning in Derr et al.~\cite{derr2018} --- that
indicates whether the surrounding case law currently supports, contests, or has weakened a
given proposition. This score has been sanity-checked against known landmark cases
(\emph{Angurbala Mullick}, \emph{Deoki Nandan v. Murlidhar}) but is explicitly not a
prediction of how a future court would rule.

\subsection{Phase III: Deterministic Verification and IRAC Synthesis}

Before any natural-language answer is produced for a judgment, a verification layer checks
every proposed AuthorityCited object against the citation ground truth: does the cited
case actually appear among this judgment's own outbound case-citation edges in the
provided metadata? This is a direct, checkable proxy for the fabricated-authority failure
mode. Matching is done on normalized party names rather than reporter citations, since the
provided metadata does not carry reporter volumes and grounding sometimes truncates
multi-citation list sentences (dropping the first party's name, e.g.\ ``Krishna Shetti v.
Gilbert Pinto'' surviving only as ``Gilbert Pinto''); either party name matching is
accepted as evidence, with a minimum length floor to avoid one-word false positives. Each
proposed authority is classified into exactly one of four states:

\begin{itemize}
\item \textsc{Verified} --- matches a real outbound case-citation edge of the judgment.
\item \textsc{Unverified} --- does not match any real edge. We further split this by
whether the proposed name appears anywhere in the judgment's own text at all, since only
the ``does not appear at all'' subset is a genuine fabrication signal; the metadata itself
is independently scraped and incomplete (it is derived from Indian Kanoon's own outbound
links, which systematically omit English and Privy Council authorities), so a real
citation the extraction garbled, or a real citation the metadata simply lacks, would also
show up as \textsc{Unverified} without this split.
\item \textsc{Self\_reference} --- the judgment citing itself or its own lower-court
history.
\item \textsc{Rejected\_not\_a\_case} --- the proposed span is not a case citation at all
(most commonly a statute section misclassified upstream; see Section~\ref{sec:results}).
\end{itemize}

Only claims that pass verification are rendered into the final IRAC-style (Issue, Rule,
Application, Conclusion) output; authorities that fail are flagged rather than silently
dropped or smoothed over into fluent prose. The verifier's heuristics were themselves
tuned against a hand audit after early mistakes: Indian reporter citations can sit outside
the parenthetical span a naive regex expects (e.g.\ ``Ayyappa Reddy, (1913) I.L.R. 38 Mad.
738''), a case-insensitive reporter pattern can accept any parenthetical including
statutory sub-clauses (``clause (iv)''), and whole-string fuzzy matching can let a shared
generic prefix carry a real mismatch (``Commissioner of Income-tax, Madras'' scoring high
against ``\ldots, Bihar and Orissa''). Fixing these on an earlier 1950s-only evaluation
run raised its verified rate from 17\% to 55\%; we treat the verifier itself as a
component requiring its own validation, not as ground truth by construction.

\section{Experimental Setup}
\label{sec:setup}

All three systems compared in this paper use the same generator model,
\texttt{llama4-scout-17b}, served from ASU's Voyager endpoint (an internal,
OpenAI-API-compatible deployment of open-weight models), and the same task: produce an
IRAC analysis of a judgment and list the prior decisions it relies on.

\textbf{Closed-book LLM.} The model receives only the case title and year --- the ``ask a
chatbot about case $X$'' setting profiled by Dahl et al.~\cite{dahl2024} --- and answers
from parametric memory alone.

\textbf{Dense RAG.} The judgment is split into approximately 1,200-character chunks,
embedded with \texttt{qwen3-embedding-8b}, and the model answers from the top-8 chunks by
cosine similarity to a fixed query about precedents, rule, and holding.

\textbf{LexGraph (ours).} Phase I grounds the first 30 keyword-matched sentences of each
judgment into the ontology, as described in Section~\ref{sec:method}.

Every authority any system lists is then scored by the identical deterministic verifier
described in Section~\ref{sec:method} (\texttt{eval\_lib.score\_authorities}), so
differences between systems reflect what they surfaced, not how leniently they were
graded. The evaluation set is 48 judgments, six drawn from each decade from the 1950s
through the 2020s via \texttt{run\_evaluation\_batch.py --per-decade 6}, from the
stratified 500-case sample described in Section~\ref{sec:dataset}. Sampling by decade
rather than by row order matters: the underlying sample file is sorted by decade, so a
naive first-$N$ selection would silently mean ``the first $N$ 1950s cases.''

We report four quantities per system: \emph{case citations}, the number of listed
authorities that are actually case references (excluding statutes, bare party names, or
other non-citation spans); \emph{verified}, how many of those match a real outbound edge
of the judgment; \emph{not named in judgment}, how many unverified citations have no party
name occurring anywhere in the judgment text at all (the sharper fabrication signal, per
Section~\ref{sec:method}); and \emph{edge recall}, the fraction of the 2,123 real outbound
case-citation edges present across these 48 judgments that a system surfaced and verified.

\section{Results}
\label{sec:results}

\subsection{Headline Comparison}

Table~\ref{tab:headline} summarizes the three systems on the 48-case evaluation set.

\begin{table*}[t]
\caption{System comparison on 48 judgments (six per decade, 1950s--2020s), scored by an identical verifier}
\label{tab:headline}
\centering
\begin{tabular}{@{}lrrrrr@{}}
\toprule
\textbf{System} & \textbf{Authorities output} & \textbf{Case citations} & \textbf{Verified} & \textbf{Not named in judgment} & \textbf{Edge recall} \\
\midrule
Closed-book LLM & 181 & 174 & 12 (6.9\%) & 101 & 11 / 2123 (0.5\%) \\
Dense RAG       & 464 & 417 & 216 (51.8\%) & 5 & 211 / 2123 (9.9\%) \\
LexGraph (ours) & 235 & 50  & 30 (60.0\%) & 0 & 30 / 2123 (1.4\%) \\
\bottomrule
\end{tabular}
\end{table*}

Closed-book prompting fabricates: only 12 of the vanilla model's 174 case citations
(6.9\%) match a real citation of the judgment, and 101 of its 162 unverified citations
name no party found anywhere in the judgment's text --- fluent, correctly formatted
citations, complete with reporter volumes, for authorities the judgment never cites. This
reproduces the failure mode Dahl et al.~\cite{dahl2024} report, on our own corpus.
Grounding in the source text, whether via RAG or via ontology-directed extraction, largely
removes it: only 5 of RAG's 201 unverified citations and 0 of LexGraph's 20 unverified
citations are absent from the judgment (Table~\ref{tab:unverified}).

\subsection{Provenance of Unverified Citations}

\begin{table}[t]
\caption{Breakdown of unverified citations by whether the proposed party name appears in the judgment text}
\label{tab:unverified}
\centering
\footnotesize
\begin{tabular}{@{}lrrrr@{}}
\toprule
\textbf{System} & \textbf{Unverified} & \textbf{Named in text} & \textbf{Not in text} & \textbf{Name too short} \\
\midrule
Closed-book & 162 & 59  & 101 & 2  \\
Dense RAG   & 201 & 185 & 5   & 11 \\
LexGraph    & 20  & 13  & 0   & 7  \\
\bottomrule
\end{tabular}
\end{table}

The remaining ``name too short'' category covers party names too short or generic (e.g.\
a single common surname) for the length floor described in Section~\ref{sec:method} to
check safely, since short-name containment matching produced false verifications during
the verifier's own hand audit (Section~\ref{sec:method}); these are left unresolved rather
than guessed in either direction. This split matters because an unverified citation is not
automatically a fabricated one.
The citation metadata used as ground truth is itself independently scraped from Indian
Kanoon's own outbound links, which systematically omit English and Privy Council
authorities, and is not guaranteed complete. A citation named in the judgment's own text
but not matched to a ground-truth edge is more likely a metadata gap, a name-matching
miss, or (for RAG and LexGraph) a real citation the extraction produced in a slightly
garbled form, than an invented authority. Under this reading, LexGraph's 0 not-in-text
unverified citations and RAG's 5 both look like genuinely low fabrication rates; the
closed-book model's 101 do not have this excuse, since nothing in a title-and-year-only
prompt could have surfaced a name that is absent from the judgment altogether.

\subsection{Per-Decade Breakdown and the Recall Bottleneck}

Table~\ref{tab:decade} shows LexGraph's verified rate and edge recall by decade.

\begin{table*}[t]
\caption{LexGraph performance by decade}
\label{tab:decade}
\centering
\begin{tabular}{@{}lrrrrrr@{}}
\toprule
\textbf{Decade} & \textbf{Cases} & \textbf{Authorities output} & \textbf{Case citations} & \textbf{Verified rate} & \textbf{Real edges} & \textbf{Edge recall} \\
\midrule
1950s & 6 & 33 & 16 & 68.8\% & 86  & 12.8\% \\
1960s & 6 & 38 & 18 & 61.1\% & 83  & 13.3\% \\
1970s & 6 & 30 & 5  & 60.0\% & 302 & 1.0\%  \\
1980s & 6 & 34 & 6  & 33.3\% & 128 & 1.6\%  \\
1990s & 6 & 21 & 2  & 0.0\%  & 196 & 0.0\%  \\
2000s & 6 & 31 & 1  & 100.0\% & 149 & 0.7\%  \\
2010s & 6 & 27 & 2  & 100.0\% & 550 & 0.4\%  \\
2020s & 6 & 21 & 0  & --      & 629 & 0.0\%  \\
\bottomrule
\end{tabular}
\end{table*}

Edge recall is concentrated almost entirely in the 1950s and 1960s (12.8\% and 13.3\%),
and falls to near zero from the 1970s onward, even though the number of real outbound
edges per decade generally grows over time (from 86 in the 1950s to 629 in the 2020s,
reflecting a more heavily cross-cited modern jurisprudence). This is consistent with the
mechanism identified in Section~\ref{sec:method}: 1950s--60s judgments announce their
cited authorities in a HEADNOTE section near the top of the document, which falls
comfortably inside a first-30-keyword-sentence window; later judgments cite authorities
throughout the reasoning, where the same fixed window captures only a small, effectively
arbitrary fraction of them. High per-decade verified rates in some later decades (100\%
in the 2000s and 2010s) are not evidence the pipeline works better there --- they come
from only 1--2 case-shaped citations being proposed at all in six judgments, i.e.\ a
precision statistic computed over almost no recall.

\subsection{Extraction Volume and Precision}

Across the 48 evaluated judgments, Phase I produced 192 Party, 362 Claim, 453 Holding, 108
Outcome, and 235 AuthorityCited objects. Of those 235 AuthorityCited objects, 180 (76.6\%)
were rejected by the verifier as not case citations at all --- almost all statute or
section spans. This is the extraction-side counterpart of the RulingAction+Patient
ambiguity noted in Section~\ref{sec:method}: the same syntactic pattern (``held that
\ldots'', ``relying on \ldots'') introduces both statutory rules and case authorities, and
the current rule set does not yet look at the content of the span (does it name a
reported case versus a numbered section) to tell them apart.

\section{Discussion and Limitations}
\label{sec:limitations}

The central, honest result of this evaluation is that \emph{plain RAG currently beats
LexGraph on coverage}. RAG recovers and verifies 9.9\% of the real citation edges in these
48 judgments against LexGraph's 1.4\%, at broadly comparable precision (51.8\% vs.\ 60.0\%
of case-shaped citations verified). The bottleneck is Phase I's candidate selection, not
Phase III's verification: LexGraph only ever grounds the first 30 keyword-matched
sentences of a judgment, so its coverage is bounded above by how much of a judgment's
citation content fits in that window, and Section~\ref{sec:results} shows that window
works only for a law-reporting convention that stops appearing after 2000. Because the
verifier is agnostic to how a candidate authority was produced, this is a solvable
integration problem rather than a limitation of the verify-before-answer approach itself:
feeding RAG-style retrieval into ontology grounding, instead of choosing between them,
should let LexGraph inherit RAG's coverage while keeping typed, auditable objects and
Phase III certification.

A second limitation is Phase I precision: 76.6\% of extracted AuthorityCited objects are
not case citations, driven by an ontology role (RulingAction+Patient) that is genuinely
overloaded across four different target concepts. More gold examples have not resolved
this in past iterations; it needs content-based disambiguation of the span itself.

A third limitation concerns the evaluation's statistical power and the ground truth it
verifies against. Six cases per decade is too small a sample to support strong per-decade
conclusions, particularly for decades where only 0--2 case-shaped citations were proposed
at all. The citation ground truth itself is Indian Kanoon's own scraped outbound-link
data, not a hand-curated authority; it is known to omit English and Privy Council
citations systematically, and 22 verified matches across the three systems in this
evaluation matched more than one candidate edge (e.g.\ a bare ``State of Orissa'') and are
flagged as ambiguous rather than confidently attributed. None of this invalidates the
comparison between systems, since all three are scored against the same imperfect ground
truth with the same matching logic, but it means absolute rates (e.g.\ ``60\% verified'')
should be read as measured against available, incomplete evidence rather than as an exact
fabrication rate.

\section{Future Work}
\label{sec:future}

The immediate next step is to replace Phase I's first-30-keyword-sentence heuristic with
retrieval-selected passages --- the same retrieval mechanism already implemented for the
RAG baseline --- as the candidate pool fed into ontology grounding, and to re-run the
full evaluation without reusing cached extractions. We also plan to run an explicit ``RAG
plus verifier'' ablation (RAG's retrieved passages, scored by Phase III without any
ontology grounding at all) to isolate how much of LexGraph's precision advantage over
plain RAG comes from typed extraction versus from the verifier itself. On the precision
side, disambiguating RulingAction+Patient by inspecting whether its filler names a
reported case or a statute section is the highest-priority ontology fix. Beyond this
project's initial scope, we intend to build a lightweight interactive dashboard so a user
can enter a land-dispute fact pattern and see the retrieved authorities, the issue-rule
-evidence trace, the Proof-Support Score, and any authorities the verifier flagged as
unsupported, which is the artifact intended for the course's live demo session.

\section{Conclusion}

Deterministic, ground-truth-checked verification works: across all three systems in this
evaluation, the ones that ground their answers in the judgment's own text (RAG and
LexGraph) show a near-zero rate of citing authorities that do not appear in that text at
all, while a closed-book LLM fabricates such citations for the majority of its unverified
output, reproducing a known legal-hallucination failure mode on this corpus. But
verification alone does not fix coverage: LexGraph's advantage in precision currently
comes at the cost of an order-of-magnitude deficit in recall relative to plain retrieval,
traced here to a specific, fixable candidate-selection bottleneck rather than to any
weakness in the verify-before-answer principle itself. The next iteration of this
pipeline combines retrieval for coverage with ontology-directed extraction for typed,
auditable objects and deterministic verification for certification, so that each stage's
contribution can be measured rather than assumed.

\section*{Acknowledgment}

We thank Prof.\ Hasan Davulcu for the project topic and dataset, and TA Anshul Trivedi for
guidance during office hours. LLM inference for this project was run on ASU's Voyager
research computing endpoint.

\begin{thebibliography}{9}

\bibitem{dahl2024}
M. Dahl, V. Magesh, M. Suzgun, and D. E. Ho, ``Large Legal Fictions: Profiling Legal
Hallucinations in Large Language Models,'' \emph{Journal of Legal Analysis}, vol. 16,
no. 1, pp. 64--93, 2024.

\bibitem{derr2018}
T. Derr, Y. Ma, and J. Tang, ``Signed Graph Convolutional Network,'' in \emph{Proc. 2018
IEEE International Conference on Data Mining (ICDM)}, 2018, pp. 929--934.

\bibitem{hoekstra2007}
R. Hoekstra, J. Breuker, M. Di Bello, and A. Boer, ``The LKIF Core Ontology of Basic Legal
Concepts,'' in \emph{Proc. 2nd Workshop on Legal Ontologies and Artificial Intelligence
Techniques (LOAIT)}, 2007.

\bibitem{legalruleml}
OASIS LegalRuleML Technical Committee, ``LegalRuleML Core Specification Version 1.0,''
\url{https://docs.oasis-open.org/legalruleml/legalruleml-core-spec/v1.0/}.

\bibitem{lewis2020}
P. Lewis \emph{et al.}, ``Retrieval-Augmented Generation for Knowledge-Intensive NLP
Tasks,'' in \emph{Advances in Neural Information Processing Systems (NeurIPS)}, vol. 33,
2020, pp. 9459--9474.

\bibitem{riloff1993}
E. Riloff, ``Automatically Constructing a Dictionary for Information Extraction Tasks,''
in \emph{Proc. Eleventh National Conference on Artificial Intelligence (AAAI-93)}, 1993.

\bibitem{sunny2026}
A. D. Sunny and I. Sivan-Sevilla, ``A Neuro-Symbolic Framework for Legal Accountability in
Public-Sector AI,'' in \emph{Proc. 2026 ACM Conference on Fairness, Accountability, and
Transparency (FAccT '26)}, 2026.

\end{thebibliography}

\end{document}
